\documentclass[a4paper, fontsize=9pt, landscape]{scrartcl}

% ============================================================
% PRÄAMBEL & LAYOUT (Standard Herr Dayi / Fachleiter Dr. Uhtes)
% ============================================================
\usepackage[a4paper, landscape, top=0.65cm, bottom=0.6cm, left=1.0cm, right=1.0cm]{geometry}
\usepackage[ngerman]{babel}
\usepackage{fontspec}
\setmainfont{Libertinus Serif}
\setsansfont{Libertinus Sans}
\usepackage{xcolor}
\usepackage{colortbl}
\usepackage{tabularx}
\usepackage{array}
\usepackage{enumitem}

\pagestyle{empty}

% Farbdefinitionen
\definecolor{primaryBlue}{RGB}{20, 50, 100}
\definecolor{lightBlue}{RGB}{245, 248, 253}
\definecolor{tableGray}{RGB}{240, 244, 250}

\renewcommand{\arraystretch}{1.08}

\begin{document}

% === KOPFZEILE METADATEN (EXAKT WIE IN DER VORLAGE) ===
\noindent
\begin{minipage}[b]{0.38\textwidth}
    {\normalsize\bfseries 2.2 Darstellung des geplanten Unterrichtsverlaufs}
\end{minipage}%
\hfill
\begin{minipage}[b]{0.61\textwidth}
    \flushright
    {\footnotesize\textbf{Thema:} Karl Marx: Die entfremdete Arbeit \enskip$\vert$\enskip \textbf{Kurs:} Q1 \enskip$\vert$\enskip \textbf{Fach:} Philosophie \enskip$\vert$\enskip \textbf{Lehrkraft:} Herr Dayi}
\end{minipage}

\vspace{0.06cm}
\hrule height 0.6pt
\vspace{0.08cm}

% === LERNZIELE (ÜBERSICHTLICH AUFGEGLIEDERT NACH DR. UHTES) ===
\noindent
{\footnotesize\textbf{\color{primaryBlue}Hauptlernziel (HLZ):} Die SuS \textbf{erschließen}, dass nach Karl Marx die industrielle Produktionsweise das menschliche Gattungswesen verkehrt, da die schöpferische Arbeit eine vierfache Entfremdung (vom Produkt, vom Tätigkeitsakt, vom Gattungswesen und vom Mitmenschen) bewirkt und freie Selbstverwirklichung in fremdbestimmte Zwangsarbeit transformiert.}

\vspace{0.08cm}
\noindent
{\footnotesize\textbf{\color{primaryBlue}Teillernziele (TLZ):} \textit{Die Schülerinnen und Schüler \dots}}
\vspace{0.06cm}
\begin{itemize}[leftmargin=1.5em, itemsep=1.5pt, topsep=2pt, parsep=0pt, label=\color{primaryBlue}\footnotesize\textbullet]
    \footnotesize
    \item \textbf{TLZ 1 (AFB I):} \textit{problematisieren} die Herrschaft der Maschine über den Menschen, \textbf{indem} sie an Charlie Chaplins Fließbandarbeit das Phänomen der Technisierung untersuchen und den Leitkonflikt der Fremdbestimmung formulieren.
    \item \textbf{TLZ 2 (AFB II):} \textit{rekonstruieren} Marx’ Argumentationskette zur Entfremdung, \textbf{indem} sie vier dekontextualisierte Kernaussagen kooperativ erschließen und am Board in eine kausale Hypothesenkette ordnen.
    \item \textbf{TLZ 3 (AFB II):} \textit{analysieren} die vier Dimensionen der entfremdeten Arbeit, \textbf{indem} sie die Textstellen im Originaltext lokalisieren und deren argumentative Funktion in einer Ergebnismatrix systematisieren.
    \item \textbf{TLZ 4 (AFB II):} \textit{verorten} Marx’ Kulturdiagnose im Advanced Organizer (Kultur als Trugbild), \textbf{indem} sie seine Position pointiert verdichten und begründet im Spannungsfeld zu Gehlen einordnen.
    \item \textbf{TLZ 5 (AFB III):} \textit{beurteilen} die Aktualität von Marx’ Paradoxon für die heutige Lebens- und Arbeitswelt, \textbf{indem} sie die Kernspannung in einem pointierten Meme gestalterisch zuspitzen.
\end{itemize}
\vspace{0.05cm}

% === TABELLARISCHER VERLAUFSPLAN (OHNE ZEITANGABEN - EXAKT WIE VORLAGE) ===
\scriptsize
\noindent
\begin{tabularx}{\textwidth}{|p{3.2cm}|X|p{3.0cm}|p{3.4cm}|}
\hline
\rowcolor{tableGray}
\textbf{\color{primaryBlue}Unterrichtsphase /} \newline \textbf{\color{primaryBlue}Funktion} & 
\textbf{\color{primaryBlue}Sach- und Handlungsaspekte (Lehrkraft- und Schüleraktivität)} & 
\textbf{\color{primaryBlue}Sozialform /} \newline \textbf{\color{primaryBlue}Handlungsmuster} & 
\textbf{\color{primaryBlue}Medien /} \newline \textbf{\color{primaryBlue}Material} \\
\hline

% --- PHASE 1 ---
\textbf{Einstieg \&} \newline \textbf{Problemorientierung} \newline
\tiny\color{primaryBlue} Problembegegnung, \newline kognitiver Konflikt, \newline Leitproblem & 
L. präsentiert den Filmausschnitt: Charlie Chaplin in \textit{Moderne Zeiten} (1936) am Fließband (monotones Schrauben-Anziehen, Hysterie, Einsaugen in die Zahnräder). \newline
\textbf{Beobachtungsauftrag:} „Was geschieht mit Charlie an der Maschine? Wie verändern sich Körper, Geist und Verhalten?“ \newline
\textbf{Murmelphase (Team):} SuS tauschen sich aus: Wer beherrscht hier wen: Mensch oder Maschine? \newline
\textbf{Problemfokussierung im Plenum:} SuS decken die paradoxe Verkehrung auf: Nicht der Arbeiter nutzt das Werkzeug, sondern die Maschine diktiert den Menschen und unterwirft seinen Organismus. \newline
L. fixiert Leitfrage am Board / Tafel: \textit{„Wer beherrscht hier eigentlich wen: Der Mensch die Maschine -- oder die Maschine den Menschen?“} \newline
L. visualisiert transparent die Struktur der Stunde (PPP Folie 3). & 
UG $\rightarrow$ \newline
Murmelphase (PA) $\rightarrow$ \newline
Plenum, \newline
Lehrkraftimpuls & 
Digitale Tafel \newline
(PPP Folie 2--3), \newline
Filmclip \newline
(Charlie Chaplin) \\
\hline

% --- PHASE 2 ---
\textbf{Fokussierung \&} \newline \textbf{Erschließung} \newline
\tiny\color{primaryBlue} Dekontextualisierte \newline Konfrontation, \newline kognitive Aktivierung & 
L. visualisiert den Arbeitsauftrag: Die SuS erhalten isolierte Kernstellen eines ungenannten Philosophen auf Streifen (Karten A--D). \newline
\textbf{Kleingruppenarbeit (Tandems / Trios):} \newline
-- SuS paraphrasieren die Kernstelle prägnant in eigenen Worten. \newline
-- SuS prüfen den Bezug zu Chaplins Fabrikschicksal und finden ein gegenwärtiges Alltagsbeispiel (z.\,B. Paketzentrum, Stechuhr, monotone Büro-/Bildschirmarbeit). \newline
\textit{Binnendifferenzierung:} Tandems arbeiten kooperativ; sprachlich heterogene Teams stützen sich gegenseitig bei der Bedeutungserschließung. & 
LV $\rightarrow$ \newline
Kleingruppenarbeit / \newline
Tandems (PA/GA) & 
Digitale Tafel \newline
(PPP Folie 4), \newline
Kernaussagen- \newline
Streifen A--D \\
\hline

% --- PHASE 3 ---
\textbf{Rekonstruktion \&} \newline \textbf{Hypothesenbildung} \newline
\tiny\color{primaryBlue} Schülermoderierte \newline Argumentations- \newline rekonstruktion, Modellierung & 
L. projiziert das interaktive \textbf{Kernaussagen-Board} am Beamer (Thesen durcheinandergewürfelt). \newline
\textbf{Impuls der Lehrkraft:} \textit{„Diese vier Sätze stammen aus ein und demselben Text und erzählen eine zusammenhängende Geschichte. In welcher logischen Reihenfolge müssen sie stehen? Was ist Ursache, was Folge?“} \newline
\textbf{Diskussion im Plenum:} SuS begründen die Kausalitätsbeziehungen und ordnen die Thesen am Beamer live in die Slots 1--4 an. \newline
SuS formulieren die gemeinsame visuelle \textbf{Kurshypothese} (Vom schöpferischen Gattungswesen $\rightarrow$ Herrschaft des Produkts $\rightarrow$ Zwangsarbeit $\rightarrow$ anthropologisches Paradoxon). & 
Plenum / \newline
moderiertes UG, \newline
SuS am Board & 
Digitale Tafel / \newline
Beamer, \newline
Interaktives \newline
Kernaussagen-Board \\
\hline

% --- PHASE 4 ---
\textbf{Erarbeitung \&} \newline \textbf{Textabgleich} \newline
\tiny\color{primaryBlue} Textbasierte Verifikation, \newline kriteriengeleitete \newline Analyse & 
1. \textbf{EA:} SuS erschließen den Originaltext aus Karl Marx’ \textit{Ökonomisch-philosophischen Manuskripten} (AB Seite 1), lokalisieren ihre Kernaussage im Textzusammenhang und notieren die genaue Zeilenangabe. \newline
2. \textbf{PA (im Team):} SuS überprüfen die am Beamer modellierte Kurshypothese am Fließtext und bestimmen die argumentative Funktion ihres Satzes sowie die jeweilige Dimension der Entfremdung. \newline
\textit{Binnendifferenzierung:} \newline
-- \textit{Förderung (Scaffolding):} Digitale Hilfestellung (Glossar, Begriffshilfen via iPad/QR-Code). \newline
-- \textit{Forderung (Sprinter):} Vergleich von Marx’ Naturbegriff mit Rousseau (Urzustand vs. bewusstes Arbeiten). & 
Einzelarbeit (EA) $\rightarrow$ \newline
Partnerarbeit (PA) & 
Arbeitsblatt (AB), \newline
Textmarker, \newline
iPads / Tablets \newline
(digitale Hilfe) \\
\hline

% --- PHASE 5 ---
\textbf{Sicherung \&} \newline \textbf{Systematisierung} \newline
\tiny\color{primaryBlue} Kategoriale Explikation, \newline Matrix-Strukturierung, \newline Ergebnissicherung & 
\textbf{Gemeinsame Auswertung im Plenum:} SuS präsentieren ihre Zeilenfunde und begründen die argumentative Funktion der 4 Textstellen. \newline
\textbf{Tafelsicherung am Smartboard:} L. und SuS sichern das 4-Stufen-Modell der Entfremdung: \newline
1. Satz A (Z.\,9--11): Freie Produktion \& Schönheit $\rightarrow$ \textit{Gattungswesen (Ideal)} \newline
2. Satz B (Z.\,18--20): Herrschaft des Kapitals über den Arbeiter $\rightarrow$ \textit{Produkt-Entfremdung} \newline
3. Satz C (Z.\,26--27): Verweigerung freier Energie / Kasteiung $\rightarrow$ \textit{Tätigkeits-Entfremdung} \newline
4. Satz D (Z.\,30--32): Mensch nur im Tierischen frei $\rightarrow$ \textit{Anthropologisches Paradoxon} \newline
\textbf{Festhalten:} SuS übertragen die Systematik in die Ergebnistabelle ihres Arbeitsblatts (Seite 2). & 
UG / Plenum $\rightarrow$ \newline
Einzelarbeit (EA) & 
Digitale Tafel / \newline
Smartboard, \newline
Arbeitsblatt (AB Seite 2) \\
\hline

% --- PHASE 6 ---
\textbf{Synthese \&} \newline \textbf{Kultur-Urteil} \newline
\tiny\color{primaryBlue} Kategoriale Verortung \newline im Advanced Organizer, \newline Urteilsbildung & 
\textbf{Leitfrage zur Synthese:} „Ist Kultur nach Marx Visitenkarte oder Trugbild?“ \newline
\textbf{Partnerarbeit (PA):} \newline
a) SuS verdichten Marx’ Kulturdiagnose in einem pointierten Leitsatz. \newline
b) SuS verorten Marx begründet im digitalen Advanced Organizer im Spannungsfeld zu Gehlen (Kultur als Entlastung) und Freud (seelische Belastung/Triebverzicht). \newline
\textbf{Kollaborative Auswertung im Plenum:} SuS begründen die Zuordnung: Kultur unter kapitalistischen Bedingungen ist kein neutraler Schutzraum, sondern institutionalisierte Entfremdung (Trugbild-Charakter). & 
Partnerarbeit (PA) $\rightarrow$ \newline
Plenum & 
Tablets / iPads \newline
(Digitaler Advanced \newline
Organizer), \newline
AB Aufgabe 3 \\
\hline

% --- DIDAKTISCHE RESERVE ---
\multicolumn{4}{|c|}{\cellcolor{tableGray}\textbf{Mögliches Stundenende / Didaktische Reserve}} \\
\hline
\textbf{Didaktische} \newline \textbf{Reserve} \newline
\tiny\color{primaryBlue} Kreativer Transfer / \newline formative Diagnose & 
\textbf{Option A (Kreativer Transfer):} SuS gestalten ein pointiertes Meme im Online-Generator (oder zeichnen auf der 4-Panel-Comic-Vorlage), das Marx’ Kernparadoxon (Mensch vs. Tier / Chaplin an der Maschine) pointiert zuspitzt. \newline
\textbf{Option B (Check-Out via digitale Tafel):} Individuelle 3-Schritt-Reflexion: „Ich weiß jetzt...“, „Besonders spannend fand ich heute...“, „Diese Frage nehme ich mit...“. & 
Einzelarbeit (EA) / \newline
Plenum & 
Digitale Tafel \newline
(PPP Folie 8--9), \newline
Meme-Generator, \newline
Comic-Vorlage \\
\hline

\end{tabularx}

\vspace{0.06cm}
\noindent
{\tiny\textbf{Verwendete Abkürzungen:} EA = Einzelarbeit, PA = Partnerarbeit, GA = Gruppenarbeit, UG = Unterrichtsgespräch, LV = Lehrervortrag, SuS = Schülerinnen und Schüler, L. = Lehrkraft, AB = Arbeitsblatt, PPP = Präsentation.}

\end{document}
